\chapter{Debate de artículos de RRHH en los equipos}

En esta dinámica trabajamos con distintos artículos que abordaban temas muy diversos, desde la negociación laboral hasta la importancia del networking o la definición de una estrategia vital. Aunque cada texto trataba aspectos diferentes, todos coincidían en una idea común: la carrera profesional no puede construirse únicamente desde lo inmediato o lo económico, sino que requiere una mirada más amplia que incorpore valores, propósito y buenas relaciones.

Uno de los debates más interesantes fue el del artículo Negotiating as a Woman of Color \cite{hbr-woc-negotiation}. Surgieron reflexiones sobre cómo los prejuicios pueden influir en determinados colectivos y sobre la necesidad de aprender a manejarlos de forma constructiva. Al leerlo pensé en la importancia que tienen las organizaciones a la hora de garantizar oportunidades reales de desarrollo y en cómo la diversidad, cuando se integra bien, se convierte en una fuente clara de innovación y crecimiento. Esto me llevó a valorar la necesidad de que las políticas de recursos humanos fomenten espacios donde todas las personas puedan participar y avanzar en igualdad.

El tema del networking también generó mucha conversación a partir del artículo Learn to Love Networking \cite{hbr-networking}. A menudo se percibe como algo incómodo o interesado, pero en realidad funciona mejor cuando se entiende como un intercambio genuino y no como una estrategia calculada. La profesora lo resumió con la idea de “nunca comas solo”, que refleja que las redes profesionales se construyen desde la generosidad. Me pareció especialmente relevante porque muestra que las conexiones no son un accesorio de la carrera, sino una parte esencial de cómo trabajamos hoy: compartir información, aprender de otros y formar parte de comunidades que multiplican las posibilidades.

La metáfora del experimento del marshmallow, incluida en el artículo How to Choose a Fulfilling Career \cite{atlantic-fulfilling-career}, ayudó a entender la importancia de desarrollar una visión a largo plazo. A partir de ahí surgieron reflexiones sobre cómo el trabajo debe ser, en sí mismo, una fuente de satisfacción, y no solo una vía para conseguir algo externo. Personalmente, me hizo pensar en cómo a veces me centro demasiado en la meta siguiente y olvido valorar el propio proceso. También recordé un consejo recurrente: no tener miedo a hacer cambios, incluso laterales, si ayudan a mantener la motivación o a explorar nuevas fortalezas.

Otro de los artículos que más impacto generó fue This Is Your Most Important Decision \cite{80000hours}. Descubrir que dedicaremos unas 80.000 horas de nuestra vida al trabajo es una forma muy directa de recordarnos que merece la pena dedicar tiempo a pensar hacia dónde queremos dirigirnos. Invertir esas horas en algo que tenga sentido, que aporte y que encaje con nuestros valores es una decisión que afecta tanto a nuestro bienestar como al funcionamiento de las organizaciones, que obtienen mejores resultados cuando las personas se sienten comprometidas.

Finalmente, artículos como How Will You Measure Your Life? \cite{christensen-measure-life} y Negotiating Your Next Job \cite{hbr-negotiating-next-job} reforzaron la importancia de pensar la carrera profesional como un proyecto coherente. No se trata de aceptar oportunidades sin más, sino de evaluar si encajan con lo que buscamos y con la manera en la que queremos vivir. En esa evaluación entran aspectos como las posibilidades de aprendizaje, el estilo de liderazgo, la cultura organizativa o el espacio para desarrollarse, elementos que a medio y largo plazo suelen ser más determinantes que el propio salario.

En conjunto, esta dinámica me ayudó a mirar la carrera con más perspectiva. Comprendí que el éxito profesional no depende solo de logros visibles, sino de la coherencia entre nuestras decisiones, los valores que nos guían y las relaciones que construimos en el camino. Cada artículo aportó una pieza distinta de ese puzzle, y el debate grupal permitió integrarlas en una visión más completa y personal de lo que significa construir una trayectoria profesional sostenible y con sentido.
